% A148: exact finite residue constraints on the two explicit elliptic lines.
\subsection{Local obstructions to unit multiples of the two elliptic points}

\begin{statement}{A.148}{Proved}
{Exact scalar residue exclusions on the cyclic-four elliptic cover}
Let \(R_P=\phi(P)\) and \(R_Q=\phi(Q)\) be the two
independent rational points of \(F\) in A.143.
For either point, no profinite unit multiple
\(zR_P\) or \(zR_Q\), with \(z\in\widehat\Z^\times\),
has a full Selmer lift to \(D^\circ\).
More generally, a lift with compact elliptic class \(zR_P\)
requires \(40\mid z\) in \(\widehat\Z\), while one with
class \(zR_Q\) requires \(10\mid z\).
The exact local residue conditions giving these conclusions are
\[
 \begin{array}{c|r|r|l}
 R&p&\operatorname{ord}(\overline R)&\mathcal A_p(R)\\ \hline
 R_P&17&10&\{0,5\}\\
 R_P&257&16&\{0,8\}\\
 R_Q&5&2&\{0\}\\
 R_Q&769&10&\{0,5\}.
 \end{array}
\tag{A.148.1}
\]
Here \(\mathcal A_p(R)\) is the set of scalar residues whose
multiple of \(\overline R\) lies in
\((\phi q)(\overline D(\F_p))\).
The map \(q:\overline D\to E\) is the actual map of A.143,
extended over the good integral model in A.146.

The complete finite image calculation at all good primes at most
\(3000\), followed by the exact necessary-residue propagation
specified below, gives the stronger divisibilities
\[
 \begin{aligned}
 D_P\mid z\quad\text{for }zR_P,\qquad
 D_Q\mid z\quad\text{for }zR_Q,\\
 D_P=2^4 3^4 5^4 7\cdot11\cdot13^2\cdot17\cdot19
              \cdot23\cdot29\cdot37\cdot41\cdot53\cdot73\cdot89,\\
 D_Q=2^4 3^4 5^3 7\cdot11^2\cdot13\cdot17\cdot19
              \cdot23\cdot29\cdot31\cdot37\cdot41\cdot53\cdot73.
 \end{aligned}
\tag{A.148.2}
\]
These are necessary local conditions for the stated full Selmer
lift; no sufficiency of the retained finite residues is asserted.
\end{statement}

\begin{proof}
\emph{The comparison with finite reduction.}
For a good prime \(p\), the map
\[
 \widehat\Z\longrightarrow F(\F_p),\qquad
              z\longmapsto (z\bmod m_p)\overline R,
       \qquad m_p=\operatorname{ord}(\overline R),
\tag{A.148.3}
\]
is the reduction of the continuous scalar map
\(z\mapsto zR\) on the completed Mordell--Weil group.
Every local ordinary point of a proposed Selmer lift extends
to the proper good model of \(\overline D\).
Every local cusp gives one of its four boundary points after
compactification.  Functoriality and local elliptic Kummer
injectivity therefore imply
\[
            z\bmod m_p\in\mathcal A_p(R),\qquad
 \mathcal A_p(R)=
       \{m\bmod m_p:m\overline R\in(\phi q)(\overline D(\F_p))\}.
\tag{A.148.4}
\]
This uses the actual proper image, so it includes all possible
boundary sources as well as ordinary sources.

There is a particularly short finite test for these image sets.
For \(R=\phi(P_*)\), where \(P_*\) is either displayed
point \(P,Q\), take \(A_m=[m]P_*\) in \(E(\F_p)\).
For a non-boundary point \((x,y)\in E(\F_p)\), the equations
of \(q\) recover
\[
                    T=\frac{1+x/n}{2},\qquad
                    W^2=\frac{y}{8k^2q_0}.
\tag{A.148.5}
\]
Thus it is in the image of \(q\) exactly when the right-hand
side is a square in \(\F_p\).
Substitution in the equation of \(E\) verifies the required
quartic equation, so this test is sufficient as well as necessary.
All four two-torsion points are in the image, by the exact
boundary fibres.
The other point of the fibre of \(\phi\) over
\(\phi(x,y)\) is its translate by \((0,0)_E\), namely
\[
                (x,y)\longmapsto
                       (-n^2/x,\ n^2y/x^2).
\tag{A.148.6}
\]
Its second coordinate differs by the nonzero square \((n/x)^2\).
Both preimages therefore have the same square test.
Since \(A_m\) already supplies one rational preimage of
\([m]\overline R\), (A.148.5) tests the full fibre and
not just one arbitrary choice of sheet.
It is independent of the representative of \(m\bmod m_p\)
for the same reason.

\emph{The two smallest detecting primes.}
At \(p=17\) the parameters in (A.146.7) reduce to
\[
              a=10,\quad q_0=10,\quad k=12,\quad n=7,
                    \quad \tfrac12=9\quad\text{in }\F_{17}.
\tag{A.148.7}
\]
The only affine points of the plane quartic have
\((T,W)=(0,0),(1,0),(9,0)\), and there is its one
infinity point.  Direct substitution over the seventeen
residues proves this complete list.
Its image on \(F\) is exactly \(\{O,(0,0)\}\).
The reduction of \(R_P\) is \((9,11)\), and its complete
cyclic orbit, in order from the origin, is
\[
 \begin{gathered}
 O,(9,11),(1,3),(8,7),(16,12),(0,0),\\
                    (16,5),(8,10),(1,14),(9,6).
 \end{gathered}
\tag{A.148.8}
\]
The chord formulas on \(v^2=u^3+8u\) verify this list
and its next return to the origin.  It proves the first row
of (A.148.1), and hence \(5\mid z\).

At \(p=5\), the parameters reduce to
\[
                    a=3,\quad q_0=4,\quad k=3,
                           \quad n=1,\quad\tfrac12=3.
\tag{A.148.9}
\]
The complete affine list consists of the three boundary
points and \(T=2\) with \(W\in\{1,2,3,4\}\).
Together with infinity its image on \(F\) is
\(\{O,(0,0),(4,0)\}\).
But \(\overline R_Q=(1,0)\) has order two on
\(v^2=u^3+4u\).
Only its zero multiple is in that image, proving the third
row of (A.148.1) and \(2\mid z\).
In particular these two primes alone exclude every profinite
unit on either of the two explicit elliptic lines.

\emph{The additional finite rows and the retained full sieve.}
For the second and fourth rows of (A.148.1), apply
(A.148.5) to the respective complete orbits of lengths
\(16\) and \(10\).  The retained certificate lists all
coordinates, square-test values, and the exact next return
to the origin, and verifies the image again by enumerating
the actual plane-quartic points.
It gives respectively only \(m=0,8\) at \(257\) and
only \(m=0,5\) at \(769\).
These give \(8\mid z\) and \(5\mid z\), proving the
stated \(40\) and \(10\) divisibilities.

For completeness, the larger necessary conditions in
(A.148.2) are certified by the following explicit finite
procedure, whose two comparison maps are just reduction
modulo the indicated divisors.
Use every prime \(p\le3000\) where \(2,k,q_0\) are units.
Compute \(m_p\) and every residue of \(\mathcal A_p(R)\)
by enumerating the finite plane quartic and its map to \(F\).
For each coefficient prime \(\ell\) dividing some \(m_p\),
let \(L_\ell\) be the largest \(\ell\)-power occurring,
and initially allow every residue in \(\Z/L_\ell\Z\).
Write \(\pi_d\) for each actual reduction map to \(\Z/d\Z\),
and put \(d_{p,\ell}=\ell^{v_\ell(m_p)}\).
Repeat these operations until no set changes:
\[
 \begin{aligned}
 \mathcal A_p&\longleftarrow
  \{a\in\mathcal A_p:
     \pi_{d_{p,\ell}}(a)\in\pi_{d_{p,\ell}}(\mathcal V_\ell)
                            \ \text{for all }\ell\mid m_p\},\\
 \mathcal V_\ell&\longleftarrow
  \{v\in\mathcal V_\ell:
     \pi_{d_{p,\ell}}(v)\in\pi_{d_{p,\ell}}(\mathcal A_p)\}.
 \end{aligned}
\tag{A.148.10}
\]
The first row is imposed for all prime divisors of \(m_p\);
the second row is imposed for each \(p\) divisible by that
coefficient prime in its order \(m_p\).
At each operation the residues of any scalar satisfying
(A.148.4) remain in every set: both rows simply compare its
two reductions to the same finite cyclic quotient.
This induction proves that every deletion is a necessary
consequence, without a converse claim about the retained sets.
All sets are finite and only decrease, so the process terminates.

The executed calculation uses \(846\) complete image constraints
for the two points and stabilizes in four rounds for each.
Taking the gcd of \(L_\ell\) and all retained
\(\mathcal V_\ell\)-residues gives exactly the prime powers
displayed in (A.148.2).  The certificate retains every deletion,
its round, its physical prime, and its actual residue modulus.
As additional exact checks, its final small-primary sets include
\[
 \begin{array}{c|c|l|c|l}
 &\multicolumn{2}{c|}{R_P}&\multicolumn{2}{c}{R_Q}\\
 \ell&L_\ell&\mathcal V_\ell&L_\ell&\mathcal V_\ell\\ \hline
 2&128&\{0,32,48,64\}&128&\{0,64,80,96\}\\
 3&243&\{0,162\}&243&\{0,81\}\\
 5&625&\{0\}&125&\{0\}\\
 7&49&\{0,21\}&49&\{0,7,14\}\\
 11&121&\{0,11,44,99\}&121&\{0\}\\
 13&169&\{0\}&169&\{0,13,156\}.
 \end{array}
\tag{A.148.11}
\]
This proves the finite strengthened divisibilities.

The exact limit of this finite conclusion is also explicit.
Let \(\mathcal T\) be the finite set of coefficient primes
occurring in the tested orders.  A scalar with
\(z_\ell=0\) on \(\mathcal T\) and \(z_\ell=1\)
at every other coefficient prime satisfies every tested
congruence, since its reduction modulo each \(m_p\) is zero.
Its exceptional set for A.146 contains every sufficiently
large odd prime.  Thus the finite calculation alone does not
place all its surviving scalars in an A.146 subgroup.
No local point or section is constructed by this surviving
finite scalar: higher residue tests and the full lifting
condition remain to be decided.
\end{proof}

The exact sources of the finite data are
\path{code/harmonic_unit_multiple_search.py},
\path{code/harmonic_scalar_residue_sieve.py}, and
\path{code/harmonic_unit_multiple_certificate.py}, with
their retained outputs and JSON data.  The first enumerates
actual finite points, the second performs the proved necessary
residue operations, and the third independently checks all
four displayed detecting rows and their finite orbits.
